\Needspace{8\baselineskip}\originalchapter{流动自由能与集体不稳定}\label{ch:instability}
\originalsection{冷束的 Doppler 平移}
均匀电子束速度 $V_0$, 线性连续性为 $-\ii(\omega-kV_0)n_1+\ii kn_0v_1=0$, 动量为 $-\ii(\omega-kV_0)mv_1=qE_1$. 联 Poisson 得 $(\omega-kV_0)^2=\omega_p^2$. 快慢空间电荷波是同一随束系振荡的 Doppler 两支. 单个冷束加均匀中性背景并不自行不稳定. 当再有另一可响应组分, 两支可能耦合并形成复根.

对静止冷离子和电子漂移, $1-\omega_{pe}^2/(\omega-kV_0)^2-\omega_{pi}^2/\omega^2=0$. 电场能与束流动能交换, 某些 k 的根有正虚部. 无碰撞增长不违反能量守恒: 被消耗的是平衡分布的流动自由能. 线性增长不能持续无限, 粒子俘获和分布重塑使线性化最终失效.

\originalsection{可解的 Lorentz 漂移模型}
取电子中心 V、离子中心0, 宽度 $v_e,v_i$, k>0:
\begin{equation}
 \epsilon=1-\frac{\omega_{pe}^2}{(\omega-kV+\ii kv_e)^2}
 -\frac{\omega_{pi}^2}{(\omega+\ii kv_i)^2}.
\end{equation}
在准中性长波低频近似忽略1, 令 $\delta=\omega_{pi}/\omega_{pe}=\sqrt{m_e/m_i}$. 两分母满足 $\omega+\ii kv_i=s\ii\delta(\omega-kV+\ii kv_e)$, s=\ensuremath{\pm}1. 解得
\begin{equation}
 \frac\omega k=\frac{\delta^2V+s\delta(v_i-v_e)}{1+\delta^2}
 -\ii\frac{s\delta V+v_i+\delta^2v_e}{1+\delta^2}.
\end{equation}
取 s 与 V 异号, 增长条件为 $|V|>v_i/\delta+\delta v_e$. 冷离子 $v_i\to0$ 时阈值为 $\delta v_e$. 准中性指电荷扰动 $en_{1i}-en_{1e}\simeq0$, 不是两个粒子数扰动之和零. 这里仍把宽度视模型参数, 不将 Lorentz 分布虚构为有限温度.

\originalsection{尺度和边界决定能否看到增长}
无限均匀空间的复频率结果要同有限装置的传播、损耗时间比较. $\gamma>0$ 但增益长度超过装置时未必观测到大幅模式; 局部正增长也不等于全局绝对不稳定. 求根时先固定 Fourier 约定、实 k 或实频率条件, 再判断根, 避免把空间衰减支误认时间增长.
